\documentclass[11pt]{article}
\usepackage{wsi-style-lua}
\usepackage{amsmath}
\newcommand{\Nesc}{\mathcal{N}_{\mathrm{esc}}}
\newcommand{\lbar}{\bar\lambda_{C}}
\wsirunhead{Bekenstein's Bound at the Compton Scale}
\title{Bekenstein's Bound at the Compton Scale of a Massive\\[2pt]
\large Elementary Particle: A Hilbert-Space Ceiling and a Numerical\\
Coincidence with the Exceptional Lie Group Sequence}
\author{Grant Whitmer}
\date{May 2026}

\begin{document}
\maketitle

\begin{abstract}
We record an elementary observation. Bekenstein's bound, evaluated at the reduced Compton
wavelength $\lbar = \hbar/(mc)$ of a massive elementary particle of rest energy $E = mc^{2}$,
gives a mass-independent saturation of $2\pi$ nats---a formal ceiling
$D \le e^{2\pi} \approx 535.49$ on the particle's internal Hilbert-space dimension. The bound is
applied at the boundary of its domain, where localization to $\lbar$ meets the pair-production
threshold; all massive Standard Model particles satisfy it comfortably ($D \le 12$). We then note
a numerical coincidence: the adjoint dimension $2\,\dim(\mathrm{adj}\,E_{8}) = 496$ sits at
$92.6\%$ of $e^{2\pi}$ ($98.8\%$ in $\log_{2}$), and the exceptional Lie sequence
$G_{2}, F_{4}, E_{6}, E_{7}, E_{8}$---which the Cartan classification terminates at $E_{8}$---climbs
monotonically toward this point. We are explicit about what this does \emph{not} show: we cannot
distinguish coincidence, selection effect, or structural significance, and because the gauge bosons
of unbroken symmetries are massless with no Compton scale, the coincidence reading is the most
defensible. We record the pattern for future reference as a short empirical observation, not a
derived constraint.
\end{abstract}

\section{The bound at the Compton scale}

The Bekenstein bound~[1] states that the entropy of any weakly-gravitating system of energy $E$
confined to a region of size $R$ satisfies
\begin{equation}
  S \le \frac{2\pi k_{B} R E}{\hbar c}.
\end{equation}
Evaluated at the reduced Compton wavelength $R = \lbar \equiv \hbar/(mc)$ of an elementary particle
with rest energy $E = mc^{2}$,
\begin{equation}
  S_{\max}(\lbar) = 2\pi k_{B}
  \quad\Longleftrightarrow\quad
  D \le \exp(2\pi) \approx 535.49,
\end{equation}
where $D$ is the dimension of the particle's internal Hilbert space (spin states $\times$ color
charge $\times$ any other gauge index). The value is mass-independent: the factor $m$ in $E$ cancels
the factor $1/m$ in $R$. This is direct cancellation in the bound, not a derivation. We denote the
function value with the escrow-series notation~[9], $\Nesc(E, L) \equiv 2\pi E L/(\hbar c)$,
recognizing that the function is Bekenstein's; we do not invoke the escrow recipe, which applies only
to the gravitational regimes of~[9] and not to a free elementary particle in vacuum.

\subsection{Domain of applicability}

Two caveats are immediate and important.

\textit{Localization at the boundary.} Confining a single elementary particle to its own reduced
Compton wavelength saturates the relativistic localization theorem~[4,5]: $\Delta x \sim \lbar$
forces $\Delta p \sim mc$, which is the threshold for pair production. A wavepacket of Compton width
is not strictly a single-particle state; particle--antiparticle pair contamination is unavoidable at
this scale. The bound~(2) is therefore being applied at the boundary of Bekenstein's formal domain.
What is well-defined in this limit is the per-species internal Hilbert dimension $D$ at the level of
the field's representation content; the spatial localization is singular and the bound is to be read
as a formal counting statement rather than a literal claim about a sharply localized particle.

\textit{Massless particles do not have a Compton scale.} The reduced Compton wavelength
$\lbar = \hbar/(mc)$ diverges as $m \to 0$. The bound~(2) therefore does not apply at any finite
scale to massless gauge bosons (photon, gluon, gauge bosons of any unbroken gauge symmetry). This
excludes the most familiar case where ``internal gauge structure'' provides large multiplicities.
These two caveats restrict the bound's reach to massive elementary particles, evaluated formally at
their own Compton scale with the understanding that the localization is at the limit of
single-particle description.

\section{Standard Model survey}

Restricted to massive elementary particles:
\begin{center}
\begin{tabular}{lcc}
\toprule
Particle & $D$ & qubits ($\log_{2} D$) \\
\midrule
Higgs scalar (physical) & 1 & 0 \\
Charged lepton (Dirac) & 4 & 2.00 \\
Quark (Dirac $\times$ 3-color) & 12 & 3.58 \\
$W^{\pm}, Z$ (3 polarizations) & 3 & 1.58 \\
\bottomrule
\end{tabular}
\end{center}
All massive Standard Model elementary particles satisfy~(2) comfortably ($D \le 12$,
$\log_{2} D \le 3.58$, against a ceiling of $9.06$ qubits). The photon and gluon are massless and
outside the bound's domain. The bound therefore imposes no constraint on observed particles. It is a
formal ceiling that becomes interesting only for hypothetical heavy particles with rich internal
structure.

\section{A numerical coincidence: the exceptional Lie group sequence}

\textit{Caveat first.} The comparison displayed in this section runs across a domain mismatch that we
acknowledge explicitly here, rather than relegating to the limitations section. The bound~(2) applies
to massive elementary particles with their own Compton wavelength. The quantity
$2\,\dim(\mathrm{adj}\,G)$ that we compare to $e^{2\pi}$ is the natural one-particle-state count for a
gauge boson of an \emph{unbroken} gauge symmetry---where the gauge boson is massless and has no finite
$\lbar$ at which to apply the bound. The two numbers therefore refer to physically distinct settings.
We display the comparison anyway because the numerical pattern is the paper's empirical content, and
the reader should see it together with the caveat rather than have either suppressed (see also
\S6(d)).

The Cartan classification of simple Lie algebras includes five exceptional algebras of finite
dimension: $G_{2}, F_{4}, E_{6}, E_{7}, E_{8}$, with adjoint dimensions $14, 52, 78, 133, 248$
respectively. Forming the one-particle-state count $D = 2\,\dim(\mathrm{adj}\,G)$:
\begin{center}
\begin{tabular}{lcccc}
\toprule
Group & $\dim\mathrm{adj}$ & $D = 2\dim(\mathrm{adj})$ & $D/e^{2\pi}$ & $\log_{2} D/\log_{2} e^{2\pi}$ \\
\midrule
$G_{2}$ & 14  & 28  & 5.2\%  & 53.0\% \\
$F_{4}$ & 52  & 104 & 19.4\% & 73.9\% \\
$E_{6}$ & 78  & 156 & 29.1\% & 80.4\% \\
$E_{7}$ & 133 & 266 & 49.7\% & 88.9\% \\
$E_{8}$ & 248 & 496 & 92.6\% & 98.8\% \\
\bottomrule
\end{tabular}
\end{center}

\textit{Metric note.} The two ratio columns describe the same underlying agreement: by definition
$\log_{2} D/\log_{2} e^{2\pi} = (\ln 2 \cdot \log_{2} D)/(2\pi)$, so the $\log_{2}$ ratio is
determined by the linear ratio. We report both because the choice of metric substantially changes the
perceived strength of the agreement ($92.6\%$ vs.\ $98.8\%$ for $E_{8}$). Reporting only the
$\log_{2}$ figure would invite an accusation of metric cherry-picking; reporting only the linear
figure would understate the visual character of the pattern.

The Cartan classification of exceptional simple Lie algebras terminates with $E_{8}$; no $E_{9}$
exists. The largest natural exceptional dimension falls below the formal Compton--Bekenstein ceiling
of~(2) by $7.4\%$ in linear units, or $1.2\%$ in $\log_{2}$ units. This is a numerical coincidence
between a transcendental constant ($e^{2\pi}$, from a localized-particle Bekenstein evaluation in the
massive-particle setting) and a Lie-algebra-classification fact ($\dim E_{8} = 248$, from the geometry
of root systems in the unbroken gauge-symmetry setting). The reader may judge its significance.

\section{Three interpretations}

In the spirit of disciplined hedging, the coincidence admits three readings.
\begin{itemize}
\item[\textbf{(A)}] \textbf{Coincidence.} $e^{2\pi}$ and $2\cdot 248$ happen to be close. The
exceptional Lie algebra classification has structural reasons (terminating at $E_{8}$ due to the
geometry of root systems and the behavior of the Weyl group) that have no obvious relation to
relativistic information bounds. \emph{Most defensible reading absent further argument.}
\item[\textbf{(B)}] \textbf{Selection effect.} Theorists prefer ``natural'' gauge groups with rich
representation theory, biasing the studied catalog toward exceptional Lie algebras. That the catalog
tops out at $E_{8}$ reflects how Lie algebras are classified, not how nature is constrained. The fit
to within $1.2\%$ may be an artifact.
\item[\textbf{(C)}] \textbf{Structural.} The Compton-scale Bekenstein number is a real physical
comparison point, and the Cartan exceptional sequence is the algebraic structure nature uses for
elementary gauge symmetries broken to the observed Standard Model. This reading requires substantive
additional input: a derivation connecting $\lbar$ to Lie-algebra dimensionality in some specific
construction.
\end{itemize}

We claim only~(A) is well-supported by the present analysis. The quantity $2\,\dim(\mathrm{adj}\,G)$
is naturally a one-particle Hilbert dimension only in the unbroken phase of the gauge symmetry, where
all gauge bosons of the multiplet share gauge invariance under $G$ and are physically
indistinguishable. In the Higgsed phase the gauge bosons separate into distinct massive species, each
carrying its own $D = 3$ (three polarizations); the full multiplet's $2\,\dim(\mathrm{adj}\,G)$ is
then a species count times a per-species dimension, not a single particle's internal dimension. The bound applies per-species in the broken phase, where each $D = 3$ is far below the ceiling.
Comparing $e^{2\pi}$ to $2\,\dim(\mathrm{adj}\,G)$ across the broken/unbroken phase transition is the
structural weakness of interpretation~(C).

\section{Relationship to the species bound and the escrow series}

The constraint that the number of light particle species is bounded by gravitational physics has a
substantial literature. Bekenstein treats the multi-species case in his original work~[1] and 2005
review~[2]. More recent quantum-gravity work~[6,7] sharpens this into the \emph{species bound}: with
$N_{\mathrm{species}}$ light fields, the effective scale of strong gravity shifts down to
$M_{P}/\sqrt{N}$, beyond which further species would induce strong-gravity effects at small energies.
Saturon and related constructions in Dvali's recent work~[8] push the same logic to bounded composite
objects.

These constraints differ from the present one in a specific way: they bound the \emph{count of
distinct species}, whereas the bound~(2) bounds the \emph{internal Hilbert dimension of a single
elementary particle} at its own Compton scale. The two are orthogonal in principle and can be applied
jointly:
\begin{equation}
  N_{\mathrm{species}} \le \text{(species bound, Dvali-type)},
  \qquad
  D_{\mathrm{internal}} \le e^{2\pi}\ \text{per species (Compton--Bekenstein)}.
\end{equation}
A theory with many species, each of small internal dimension, can saturate the species bound without
saturating the per-species Compton--Bekenstein bound, and conversely. The present paper addresses only
the per-species side, which has been less emphasized in the literature. The function notation $\Nesc$
is borrowed from~[9], which formalizes a cross-regime observation about the static escrow construction $|U|/T$ producing the Bekenstein-bound form in three gravitational regimes (test mass at horizon, black-hole horizon, Rindler-wedge perturbation). The recipe $|U|/T$ does not apply to a free elementary particle in vacuum (no gravitational binding energy, no horizon temperature), so the present work is not an extension of the gravitational entropy escrow series~[9--14]---we use the notation only.

\section{Limitations}
\begin{itemize}
\item[\textbf{(a)}] \textbf{Bound at boundary of applicability.} Eq.~(2) is applied where
single-particle localization breaks down via pair production. The bound is to be read formally, with
$D$ as the well-defined object.
\item[\textbf{(b)}] \textbf{Composite resolution.} Particles with substructure at scales
$L > \lbar$ (hadrons at the QCD scale) carry information at the larger scale, and the bound is not the
operative constraint.
\item[\textbf{(c)}] \textbf{Massless gauge bosons not constrained.} Photons, gluons, and gauge bosons
of any unbroken gauge symmetry have no finite $\lbar$; they are outside the bound's domain.
\item[\textbf{(d)}] \textbf{Species count versus internal dimension.} The identification ``one gauge
boson of group $G$ has $D = 2\,\dim(\mathrm{adj}\,G)$'' in \S3 treats the full adjoint multiplet as a
single object's internal Hilbert space; in the Higgsed phase the gauge bosons separate into distinct
massive species, each with $D = 3$. The numerical coincidence stands as a fact about numbers, not a
physics constraint.
\item[\textbf{(e)}] \textbf{Not an extension of the escrow framework.} The present work uses $\Nesc$
as notation but does not invoke the static escrow recipe $|U|/T$.
\end{itemize}

\section{Conclusion}

For a massive elementary particle, Bekenstein's bound evaluated at its reduced Compton wavelength
yields a mass-independent saturation value of $2\pi$ nats, equivalently a formal Hilbert-space dimension
ceiling of $D \le e^{2\pi} \approx 535$. The bound is applied at the boundary of its formal domain
(single-particle localization at the pair-production threshold). Standard Model massive particles
satisfy the ceiling comfortably.

We record the numerical coincidence that the largest exceptional Lie algebra dimension,
$2\,\dim(\mathrm{adj}\,E_{8}) = 496$, sits $1.2\%$ below the ceiling in $\log_{2}$ space, and that the
exceptional sequence $G_{2}, F_{4}, E_{6}, E_{7}, E_{8}$ climbs monotonically toward this point. We do
not claim physical significance for this coincidence. The mismatch in domain of applicability between
the bound (massive particles) and the natural setting of $2\,\dim(\mathrm{adj}\,G)$ (unbroken gauge
symmetry, massless gauge bosons) makes the structural interpretation~(C) untenable without further
argument we do not provide. The pattern is presented for future reference as a short empirical
observation.

\subsection*{Methodological note}
This paper is a revised working draft that began with a broader claim: that the Compton-scale
Bekenstein bound carves the gauge-group landscape of fundamental physics, with $E_{8}$ saturating to
$98.8\%$ as evidence of structural significance. Peer review identified five technical problems with
that framing: (i) the corollary positioning within the escrow series was undercut by the
recipe-versus-function distinction; (ii) single-particle localization at $\lbar$ contends with pair
production; (iii) $D = 2\,\dim(\mathrm{adj}\,G)$ conflates a per-species one-particle-state count with
the broken-phase massive-species count; (iv) the heterotic-string preference argument applied the
bound to massless gauge bosons of extended objects; (v) the comparison requires the bound to apply at
the Compton scale of objects that have no such scale. The present version is what survives these
objections. The pattern that motivated the earlier draft remains intriguing enough to record; the
framework claims around it did not survive.

\subsection*{Data, Code \& Resources}
\noindent Manuscript source and archived record: \url{https://github.com/Windstorm-Institute/compton-corollary} (Zenodo: \href{https://doi.org/10.5281/zenodo.20163450}{10.5281/zenodo.20163450}). Experiments, datasets, and analysis code: \url{https://github.com/Windstorm-Labs/compton-corollary}. Windstorm Institute: \url{https://windstorminstitute.org}. Windstorm Labs: \url{https://windstormlabs.com}.

\begin{thebibliography}{99}\small
\bibitem{bek1} J.~D. Bekenstein, \textit{Universal upper bound on the entropy-to-energy ratio for
bounded systems}, Phys.\ Rev.\ D \textbf{23}, 287 (1981).
\bibitem{bek2} J.~D. Bekenstein, \textit{How does the entropy/information bound work?}, Found.\ Phys.\
\textbf{35}, 1805 (2005). arXiv:quant-ph/0404042.
\bibitem{casini} H.~Casini, \textit{Relative entropy and the Bekenstein bound}, Class.\ Quantum Grav.\
\textbf{25}, 205021 (2008). arXiv:0804.2182.
\bibitem{wigner} E.~P. Wigner, \textit{On unitary representations of the inhomogeneous Lorentz group},
Ann.\ Math.\ \textbf{40}, 149 (1939).
\bibitem{nw} T.~D. Newton and E.~P. Wigner, \textit{Localized states for elementary systems}, Rev.\
Mod.\ Phys.\ \textbf{21}, 400 (1949).
\bibitem{dvali1} G.~Dvali, \textit{Black holes and large $N$ species solution to the hierarchy
problem}, Fortschr.\ Phys.\ \textbf{58}, 528 (2010). arXiv:0706.2050.
\bibitem{dvali2} G.~Dvali and M.~Redi, \textit{Black hole bound on the number of species and quantum
gravity at LHC}, Phys.\ Rev.\ D \textbf{77}, 045027 (2008). arXiv:0710.4344.
\bibitem{dvali3} G.~Dvali, \textit{Entropy bound and unitarity of scattering amplitudes}, J.\ High
Energy Phys.\ \textbf{03}, 126 (2021). arXiv:2003.05546.
\bibitem{nesc} G.~L. Whitmer, \textit{The Nesc Recipe: A Cross-Regime Observation of Bekenstein-Bound
Saturation from the Static Escrow Construction $|U|/T$}, Windstorm Institute (2026).
DOI:10.5281/zenodo.20145105.
\bibitem{escrow} G.~L. Whitmer, \textit{Spacetime as Escrow Bookkeeping}, Windstorm Institute (2026).
DOI:10.5281/zenodo.20126090.
\bibitem{lattice} G.~L. Whitmer, \textit{A Lattice QFT Test of the Static Escrow Postulate},
Windstorm Institute (2026). DOI:10.5281/zenodo.20057537.
\bibitem{c8} G.~L. Whitmer, \textit{On the Status of Local Entropy-Current Extensions of the
Gravitational Entropy Escrow Framework}, Windstorm Institute (2026). DOI:10.5281/zenodo.20041991.
\bibitem{synth} G.~L. Whitmer, \textit{Gravitational Entropy Escrow: An Interpretive Synthesis},
Windstorm Institute (2026). DOI:10.5281/zenodo.20031931.
\bibitem{phonon} G.~L. Whitmer, \textit{A Non-Equilibrium Efficiency Bound for Phonon Extraction in
BEC Analog Gravity}, Windstorm Institute (2026). DOI:10.5281/zenodo.20014390.
\end{thebibliography}

\end{document}
